%!TEX root = ../main.tex
%Downwash distribution: The variation of the downward velocity component induced by a wing across its span, affecting lift and drag characteristics.
\section{Literature study}

Ground-effect vehicles have been the subject of sustained aerodynamic research since the mid-twentieth century, initially motivated by military transport applications, and later extended to civilian and commercial concepts. Despite the amount of existing studies, uncertainties remain regarding the aerodynamic behaviour of lifting surfaces operating in close proximity to the ground, particularly with respect to stability, sensitivity to operating height, and the applicability of conventional airfoil design principles. This literature study reviews the current state of knowledge on ground-effect aerodynamics, with emphasis on lift and drag characteristics, pressure distribution effects, and design implications relevant to GEV operation.\\

\subsection{Ground Effect Aerodynamics}

When a lifting surface operates in close proximity to the ground, its aerodynamic characteristics differ significantly from those observed in free-flight conditions. Ground effect is primarily characterised by an increase in lift coefficient ($C_L$) and a reduction in induced drag ($C_{D,i}$), resulting in an improved lift-to-drag ratio compared to operation far from the ground \cite{WIGE-vehicles}. These effects arise from the interaction between the wing-induced flow field and the ground plane, which alters both the downwash distribution and the pressure field around the wing.\\

\noindent In free flight, wingtip vortices generate a downward-induced velocity component that reduces the effective angle of attack across the span. As the wing approaches the ground, the formation and strength of these vortices are suppressed, leading to a reduction in downwash and a corresponding increase in effective angle of attack. This mechanism contributes directly to the observed lift enhancement and induced drag reduction in ground effect \cite{IJET_GEV_investigation}\cite{ASME_CFD_simulation}.\\

\noindent In addition to vortex-related effects, ground proximity modifies the pressure distribution on the lifting surface. The presence of the ground restricts the downward displacement of streamlines beneath the wing, resulting in elevated static pressure on the lower surface and reduced pressure recovery losses. This altered pressure field enhances lift generation while simultaneously influencing pitching moment characteristics, which can have important implications for longitudinal stability.\\

\noindent The magnitude of ground-effect influences is commonly described using the height-to-chord ratio ($h/c$) \cite{ELSEVIER_transonic_and_supersonic}. Experimental and numerical studies indicate that significant aerodynamic changes occur for $h/c \leq 0.5$, with increasingly pronounced effects as the wing approaches the ground. In so-called extreme ground-effect (EGE) conditions, typically defined as $h/c \leq 0.1$ \cite{WIGE-vehicles}, aerodynamic behaviour departs substantially from free-flight assumptions and becomes highly sensitive to small changes in height.\\

\subsection{Airfoil Performance in Ground Effect}

The aerodynamic behaviour of airfoils operating in ground effect is governed by two complementary mechanisms: chord-dominated ground effect (CDGE) and span-dominated ground effect (SDGE). These mechanisms act at different spatial scales and jointly determine the lift, drag, and stability characteristics of ground-effect vehicles.\\

\noindent Chord-dominated ground effect arises from the confinement of the flow between the wing and the ground plane. As the clearance decreases, the air beneath the wing is compressed into a narrowing channel, converting dynamic pressure into increased static pressure, often referred to as the ram effect. Experimental and numerical studies show that this process leads to elevated pressure on the lower surface and reduced flow velocities on both the pressure (lower) and suction (upper) sides of the airfoil. Although suction on the upper surface is diminished, the substantial increase in lower-surface pressure results in a net gain in lift. CDGE therefore links airfoil performance directly to sectional geometry, particularly camber distribution, thickness, and leading-edge pressure recovery.\\

\noindent Span-dominated ground effect, by contrast, is associated with the suppression of wingtip vortices as the wing approaches the ground. The ground plane interferes with the development of the vortex system, reducing downwash and effectively increasing the wing's aspect ratio. This mechanism primarily influences induced drag and becomes increasingly significant for wings with larger spans or the presence of endplates. While SDGE is a wing-level phenomenon rather than a purely sectional one, its interaction with airfoil selection remains crucial, as pitching moment characteristics and lift sensitivity determine the vehicle's stability as height varies \cite{ASME_CFD_simulation}.\\

\noindent A number of airfoil families have been investigated for their suitability in ground-effect operation. Among these, the DHMTU series, developed at the Department of Hydro-Mechanics of the Marine Technical University in St. Petersburg, was explicitly designed to exploit ground-effect conditions. Characterised by its eight-digit designation (DHMTU-a-b-c-d-e-f-g-h), the series prioritises high lift-to-drag ratios and controllability at low ground clearances. Comparative studies often evaluate DHMTU profiles against conventional airfoils such as the NACA 4412 and Clark Y. Although not specifically designed for GEVs, these airfoils possess relatively flat lower surfaces and moderate camber, which can be advantageous under CDGE-dominated conditions \cite{WIGE-vehicles}.\\

\noindent Beyond sectional geometry, research has also explored compound wing configurations that combine a rectangular centre section with tapered or reverse-tapered outer panels. Such designs seek to balance CDGE benefits at the centreline with SDGE-driven induced-drag reductions near the wingtips. Results suggest that this hybrid approach can improve aerodynamic efficiency in extreme ground-effect regimes \cite{Compound_wing_investigation}.\\

\noindent Airfoil optimisation in ground effect involves several competing geometric considerations. Increased camber thickness enables lift generation at lower angles of attack and mitigates severe leading-edge pressure peaks that promote early flow separation. Camber shape is equally important: S-shaped or reflex camber profiles are often employed to improve longitudinal stability by counteracting nose-down pitching moments as height decreases. Airfoil thickness also plays a significant role; thicker profiles (typically 12\%--15\% thickness-to-chord ratio) tend to generate higher lift at low incidence and exhibit smoother pressure gradients. However, this comes at the cost of increased profile drag at higher clearances. Conversely, thinner profiles (around 9\% thickness) have been shown to achieve improved lift-to-drag ratios at moderate heights ($0.1 < h/c < 0.3$), highlighting the trade-offs inherent in airfoil selection for ground-effect vehicles \cite{ISU_thesis}\cite{Rozhdestvensky_WIG-book}\cite{SPRINGER_WIG-book}.\\

\subsection{Stability Considerations for GEVs}

The aerodynamic performance and safety of wing-in-ground-effect (WIG) vehicles are governed by a strong coupling between pitch attitude $ \alpha $ and flight height $ h $ above the surface. Unlike conventional aircraft, whose longitudinal stability is primarily determined by the relative position of the center of gravity and a nearly fixed aerodynamic center, GEVs operate in a regime where aerodynamic forces vary rapidly with height. As a result, small vertical disturbances can induce significant changes in lift and pitching moment, making longitudinal stability a central design challenge. This is also illustrated by Irodov's criteria, which can be represented as follows:

\[X_{\alpha} - X_{H} > 0 \qquad(with\,H = h/c)\]
\[X_{\alpha} = \frac{\partial C_m}{\partial \alpha}\Big/ \frac{\partial C_L}{\partial \alpha}\]
\[X_{H} = \frac{\partial C_m}{\partial H}\Big/ \frac{\partial C_L}{\partial H}\]\\

\noindent Where $X_{\alpha}$ and $X_{H}$ are the aerodynamic centers of the angle of attack, $\alpha$, and height ratio, $H$, respectively \cite{Longitudinal_stability_study}.\\

\noindent A defining feature of GEV aerodynamics is the height-dependent migration of the center of pressure. As ground clearance decreases, the buildup of static pressure beneath the wing intensifies due to restricted flow and partial stagnation. This pressure increase is not uniform: the high-pressure region expands aft toward the trailing edge, causing the center of pressure to shift rearward. In extreme ground effect, this shift can move the center of pressure from the quarter-chord position typical of free flight toward approximately the half-chord location. The resulting change in pitching moment produces a pronounced nose-down tendency that must be counteracted to maintain trim \cite{ISU_thesis}.\\

\noindent This height sensitivity introduces an inherent form of longitudinal instability. During steady ground-effect flight, restoring forces can be maintained within a narrow altitude band. However, during transitions out of ground effect — such as climb-out or “jump” maneuvers — the sudden loss of ram-pressure support causes the center of pressure to migrate forward rapidly. This abrupt shift generates a strong nose-up pitching moment, increasing the risk of stall or loss of control if not properly managed. Consequently, many WIG craft are stable only within a limited “altitude corridor” close to the surface \cite{IFAC_new-gen_control_systems}.\\

\noindent The pitching-moment characteristics of the airfoil play a critical role in mitigating these stability challenges. Conventional airfoils typically exhibit increasingly negative pitching moments as ground clearance decreases, exacerbating trim requirements. In contrast, airfoils specifically designed for ground-effect operation often employ reflexed or S-shaped camber lines and relatively flat lower surfaces. These features reduce the sensitivity of the pitching moment to height variations by unloading the aft portion of the wing, thereby limiting center-of-pressure migration. By stabilizing the pitching-moment behavior, such airfoils reduce the dependence on oversized horizontal tailplanes, which were commonly required in early airplane-type WIG designs \cite{Conference_paper_DHMTU}\cite{DSTO}.

\subsection{Design Features Influencing Ground Effect}

The aerodynamic efficiency of Wing-In-Ground (WIG) vehicles is strongly governed by the
interaction between the lifting surfaces and the ground plane. Unlike conventional aircraft,
where aerodynamic performance is largely independent of altitude once clear of the boundary
layer, WIG craft operate within a constrained flow environment in which geometric design
choices directly influence lift generation, drag reduction, and stability.\\

\noindent A fundamental distinction between WIG craft and conventional aircraft is the systematic use
of low geometric aspect ratios (AR). While high aspect ratios are generally favored in free flight to reduce induced drag, WIG vehicles exploit the proximity of the ground to suppress
flight to reduce induced drag, WIG vehicles exploit the proximity of the ground to suppress
wingtip vortices, thereby artificially increasing the effective aspect ratio. As a result,
low-AR wings operating in ground effect can achieve lift-to-drag ratios comparable to much
larger wings in free air. In addition to aerodynamic benefits, low-AR configurations offer
greater structural stiffness and reduce sensitivity to vertical gusts, mitigating potentially
dangerous oscillatory motions near the surface \cite{IFAC_controlled_WIG_flight_concept}\cite{HINDAWI_unsteady_effects}.\\

\noindent The effectiveness of ground effect is further enhanced through the use of wingtip devices such
as endplates or sideplates, which act as sealing elements that limit the lateral leakage of
high-pressure air from beneath the wing. Experimental and numerical studies consistently show
that endplates are most effective on low-aspect-ratio wings and when the clearance between the
endplate lower edge and the ground is minimized. A significant portion of the aerodynamic
benefit arises from endplate surfaces extending below the wing's lower surface, reinforcing
the pressure build-up beneath the wing. The use of endplates or fences on lifting surfaces—
including main wings, composite wings, and horizontal stabilizers—has been reported to
increase aerodynamic efficiency by up to 25\% compared to configurations without such devices \cite{ISU_thesis}\cite{Rozhdestvensky_WIG-book}\cite{SPRINGER_WIG-book}.\\

\noindent Wing placement relative to the fuselage represents another critical design parameter. The
literature overwhelmingly favors low-wing configurations, as positioning the wing closer to
the surface allows the trailing edge to operate within the region of strongest ground effect,
typically characterized by a clearance below approximately 25\% of the chord length. In
integrated configurations such as Dynamic Air Cushion Wing-In-Ground (DACWIG) concepts, the
fuselage, wing, and lateral buoyancy elements are blended to form a continuous cavity beneath
the craft. This approach reduces structural stress concentrations, minimizes the weight
penalty of non-lifting fuselage components, and improves the confinement of the air cushion,
thereby enhancing overall aerodynamic efficiency \cite{Rozhdestvensky_WIG-book}\cite{SPRINGER_WIG-book}\cite{IFAC_controlled_WIG_flight_concept}.\\

\noindent Collectively, these design features demonstrate that successful WIG vehicle performance
depends not only on airfoil selection, but on a holistic integration of wing geometry, wingtip
devices, and fuselage layout. The strong coupling between aerodynamic efficiency and geometric
configuration underscores the importance of tailored design strategies that explicitly
account for ground-effect-specific flow phenomena.

\subsection{Research Gaps}

Despite the well-documented aerodynamic advantages of Wing-In-Ground (WIG) effect technology,
several persistent research gaps continue to limit its transition from experimental concepts to reliable operational platforms. These gaps primarily stem from the limited availability of public performance data, inconsistencies between numerical predictions and experimental observations, and an incomplete understanding of longitudinal stability and whole-vehicle integration in ground-effect flight.\\

\noindent A major impediment identified in the literature is the scarcity of detailed, publicly available
data on airfoils specifically designed for sustained operation in close ground proximity. Much of the foundational research—particularly studies related to the DHMTU airfoil series developed in the former Soviet Union—remains classified or poorly documented in open academic sources \cite{ASME_CFD_simulation}. As a result, contemporary researchers frequently rely on conventional NACA sections or reverse-engineered profiles that were originally optimized for free-flight conditions. This reliance limits the ability to accurately capture the pressure distributions and pitching moment characteristics unique to ground effect. The literature therefore emphasizes the need for a comprehensive, openly accessible airfoil database tailored to ground-effect applications.\\

\noindent Another recurring issue is the inconsistency between Computational Fluid Dynamics (CFD) results
and experimental data, particularly in extreme ground-effect regimes. While numerical methods have advanced significantly, reduced-order approaches such as the Vortex Lattice Method (VLM) often overpredict lift due to their inability to adequately represent viscous effects and airfoil thickness at very low clearances. Experimental validation is further complicated by limitations in wind tunnel testing, where fixed-floor configurations allow the development of an artificial boundary layer that distorts the under-wing flow. Accurate modeling of ground effect requires the use of rolling-floor facilities or moving-boundary conditions, which remain relatively scarce \cite{ASME_CFD_simulation}.\\

\noindent Longitudinal stability and control constitute some of the most critical unresolved challenges
in WIG vehicle design \cite{Longitudinal_stability_study}. Although theoretical frameworks such as Irodov's stability criteria provide valuable insight, much of the existing research prioritizes cruise efficiency metrics,
such as lift-to-drag ratio, over the strongly non-linear coupling between pitch angle and height that dominates near-surface flight. Furthermore, there is no clear consensus regarding optimal stabilized operating altitudes or the effectiveness of automated control systems for mitigating disturbances caused by waves, gusts, and rapid changes in ground clearance.\\

\noindent Finally, the literature reveals a lack of comparative studies across different airfoil families
and a limited emphasis on treating the WIG vehicle as an integrated system \cite{HINDAWI_Exprerimental_investigation}. Many investigations focus on isolated wing aerodynamics while neglecting hydrodynamic effects during take-off, aeroelastic deformation of lightweight structures, and the interaction between aerodynamic and hydrodynamic loads. Addressing these deficiencies will require a shift toward coupled, time-domain simulations that integrate aerodynamics, hydrodynamics, and structural dynamics to ensure stability and safety across the full operational envelope.
